% ============================================================
%  Restauro de Passé — Baralho completo (54 cartas)
%  Compile com: pdflatex cartas.tex  (duas vezes)
%  Interface da \carta: {cor}{naipe}{icon}{título}{valor}
%    {ilustração}{citação}{saber}{efeito} + \cartaB: {código}{cor peça}{inicial}
% ============================================================
% ============================================================
%  Restauro de Passé — Motor de cartas (TikZ)
%  Estilo inspirado no exemplo "Creating playing cards using TikZ"
%  (https://tex.stackexchange.com/questions/47924), recriado do zero.
% ============================================================
\documentclass[10pt]{article}

\usepackage[a4paper,margin=8mm]{geometry}
\usepackage[T1]{fontenc}
\usepackage[utf8]{inputenc}
\usepackage[portuguese]{babel}
\usepackage{lmodern}
\usepackage{microtype}
\usepackage{xcolor}
\usepackage{tikz}
\usetikzlibrary{calc,babel}
\usepackage{helvet}
\usepackage{multicol}

% ---------- Paleta (tons de terra do Recôncavo) ----------
\definecolor{fachada}{HTML}{B04A2E}   % tijoco queimado
\definecolor{nave}{HTML}{2E6E8C}      % azul da baía
\definecolor{capela}{HTML}{6C4B99}    % roxo litúrgico
\definecolor{taipa}{HTML}{7A5A33}     % barro / taipa
\definecolor{telhado}{HTML}{3F7E58}   % verde de massa + cobre
\definecolor{historia}{HTML}{2E5E8C}  % naipe História
\definecolor{construcao}{HTML}{4E7C31}% naipe Construção
\definecolor{acao}{HTML}{B03A2E}      % naipe Ação
\definecolor{desafio}{HTML}{C4841D}   % naipe Desafio
\definecolor{ouro}{HTML}{E9B44C}      % caixa de valor
\definecolor{tinta}{HTML}{333333}     % cinza do título
\definecolor{areia}{HTML}{F2EDE1}     % fundo da ilustração
\definecolor{calcario}{HTML}{FAF6EC}  % fundo de texto

% ---------- Medidas ----------
\newlength{\larguracarta}\setlength{\larguracarta}{63mm}
\newlength{\alturacarta}\setlength{\alturacarta}{88mm}

% ---------- Pictogramas (estilo xilogravura de linha) ----------
% Todos desenhados numa caixa (0,0) -- (1,1), traço escuro na cor #1.
\newcommand{\pictograma}[1]{#1}

% Igreja de fachada com porta e janela em arco
\newcommand{\icoIgreja}[1]{%
  \draw[color=#1,line width=1.1pt,line join=round]
    (0.18,0.12) -- (0.18,0.62) -- (0.5,0.82) -- (0.82,0.62) -- (0.82,0.12) -- cycle
    (0.40,0.12) -- (0.40,0.38) arc[start angle=180,end angle=0,radius=0.1] -- (0.60,0.12)
    (0.50,0.50) circle[radius=0.055]
    (0.5,0.82) -- (0.5,0.92);
  \draw[color=#1,line width=1.1pt] (0.455,0.87) -- (0.545,0.87);}

% Torre sineira
\newcommand{\icoTorre}[1]{%
  \draw[color=#1,line width=1.1pt,line join=round]
    (0.32,0.10) -- (0.32,0.68) -- (0.68,0.68) -- (0.68,0.10)
    (0.26,0.68) -- (0.5,0.86) -- (0.74,0.68)
    (0.44,0.55) -- (0.44,0.63) arc[start angle=180,end angle=0,radius=0.06] -- (0.56,0.55) -- cycle
    (0.5,0.86) -- (0.5,0.95);
  \draw[color=#1,line width=1.1pt] (0.462,0.905) -- (0.538,0.905);}

% Molde de taipa de pilão com pilão
\newcommand{\icoTaipa}[1]{%
  \draw[color=#1,line width=1.1pt,line join=round]
    (0.16,0.14) rectangle (0.84,0.42)
    (0.16,0.28) -- (0.84,0.28)
    (0.50,0.90) -- (0.50,0.44)
    (0.41,0.90) -- (0.59,0.90) -- (0.59,0.80) -- (0.41,0.80) -- cycle
    (0.20,0.10) -- (0.80,0.10);}

% Pilão de taipa caindo (tração de corrente)
\newcommand{\icoPilao}[1]{%
  \draw[color=#1,line width=1.1pt,line join=round]
    (0.20,0.90) -- (0.20,0.18) -- (0.80,0.18) -- (0.80,0.90)
    (0.34,0.90) -- (0.34,0.50) -- (0.66,0.50) -- (0.66,0.90)
    (0.40,0.56) -- (0.60,0.56) (0.40,0.62) -- (0.60,0.62) (0.40,0.68) -- (0.60,0.68);
  \draw[color=#1,line width=1.1pt] (0.20,0.90) -- (0.14,0.97) (0.80,0.90) -- (0.86,0.97);}

% Telha colonial
\newcommand{\icoTelha}[1]{%
  \draw[color=#1,line width=1.1pt,line join=round]
    (0.14,0.30) .. controls (0.30,0.62) and (0.42,0.66) .. (0.5,0.66)
    (0.86,0.30) .. controls (0.70,0.62) and (0.58,0.66) .. (0.5,0.66)
    (0.14,0.30) -- (0.24,0.30) (0.86,0.30) -- (0.76,0.30)
    (0.5,0.66) -- (0.5,0.76);
  \draw[color=#1,line width=1.1pt] (0.20,0.16) -- (0.80,0.16) (0.24,0.22) -- (0.76,0.22);}

% Barco / canoa de pescadores
\newcommand{\icoBarco}[1]{%
  \draw[color=#1,line width=1.1pt,line join=round]
    (0.14,0.34) -- (0.20,0.18) -- (0.80,0.18) -- (0.86,0.34) -- cycle
    (0.50,0.34) -- (0.50,0.72) -- (0.76,0.42) -- (0.50,0.42)
    (0.10,0.08) .. controls (0.30,0.02) and (0.40,0.12) .. (0.50,0.08)
    .. controls (0.60,0.04) and (0.72,0.12) .. (0.90,0.08);}

% Peixe
\newcommand{\icoPeixe}[1]{%
  \draw[color=#1,line width=1.1pt,line join=round]
    (0.20,0.48) .. controls (0.38,0.72) and (0.62,0.72) .. (0.72,0.48)
    .. controls (0.62,0.24) and (0.38,0.24) .. (0.20,0.48)
    (0.72,0.48) -- (0.88,0.62) -- (0.88,0.34) -- cycle;
  \draw[color=#1,line width=1.1pt,fill=#1] (0.30,0.52) circle[radius=0.022];}

% Mão / ação (martelo levantado)
\newcommand{\icoAcao}[1]{%
  \draw[color=#1,line width=1.1pt,line join=round]
    (0.30,0.14) -- (0.30,0.52) arc[start angle=180,end angle=0,radius=0.06] --
    (0.42,0.52) arc[start angle=180,end angle=0,radius=0.06] --
    (0.54,0.52) arc[start angle=180,end angle=0,radius=0.06] --
    (0.66,0.52) -- (0.66,0.30)
    (0.30,0.52) -- (0.24,0.66) -- (0.36,0.78) -- (0.60,0.78) -- (0.72,0.66) -- (0.66,0.52)
    (0.60,0.78) -- (0.66,0.88);
  \draw[color=#1,line width=1.1pt,fill=#1] (0.62,0.88) circle[radius=0.05];}

% Interrogação em círculo
\newcommand{\icoDesafio}[1]{%
  \draw[color=#1,line width=1.1pt] (0.5,0.5) circle[radius=0.34];
  \node[color=#1,font=\sffamily\bfseries\Large] at (0.5,0.47) {?};}

% Sol do Recôncavo
\newcommand{\icoSol}[1]{%
  \draw[color=#1,line width=1.1pt] (0.5,0.52) circle[radius=0.16];
  \foreach \a in {0,45,...,315}{%
    \draw[color=#1,line width=1.1pt] ($(0.5,0.52)+(\a:0.24)$) -- ($(0.5,0.52)+(\a:0.34)$);}
  \draw[color=#1,line width=1.1pt] (0.22,0.16) -- (0.78,0.16);}

% Pergaminho / registro histórico
\newcommand{\icoPergaminho}[1]{%
  \draw[color=#1,line width=1.1pt,line join=round]
    (0.26,0.18) -- (0.26,0.82) (0.74,0.18) -- (0.74,0.82)
    (0.26,0.82) .. controls (0.36,0.72) and (0.38,0.90) .. (0.5,0.80)
    .. controls (0.62,0.90) and (0.64,0.72) .. (0.74,0.82)
    (0.34,0.64) -- (0.66,0.64) (0.34,0.52) -- (0.66,0.52) (0.34,0.40) -- (0.58,0.40);}

% Arco / pedra de cantaria
\newcommand{\icoArco}[1]{%
  \draw[color=#1,line width=1.1pt,line join=round]
    (0.26,0.14) -- (0.26,0.48) arc[start angle=180,end angle=0,radius=0.24] -- (0.74,0.14)
    (0.36,0.14) -- (0.36,0.46) arc[start angle=180,end angle=0,radius=0.14] -- (0.64,0.14)
    (0.14,0.08) -- (0.86,0.08);}

% Casas do povoado (moradores)
\newcommand{\icoCasas}[1]{%
  \draw[color=#1,line width=1.1pt,line join=round]
    (0.14,0.14) -- (0.14,0.42) -- (0.32,0.54) -- (0.50,0.42) -- (0.50,0.14)
    (0.54,0.14) -- (0.54,0.50) -- (0.72,0.62) -- (0.86,0.50) -- (0.86,0.14)
    (0.23,0.14) -- (0.23,0.30) -- (0.32,0.30) -- (0.32,0.14)
    (0.63,0.14) -- (0.63,0.34) -- (0.73,0.34) -- (0.73,0.14);}

% Galeria de arcos (colunas e arcos)
\newcommand{\icoGaleria}[1]{%
  \draw[color=#1,line width=1.1pt,line join=round]
    (0.16,0.14) -- (0.16,0.86) (0.38,0.14) -- (0.38,0.86)
    (0.62,0.14) -- (0.62,0.86) (0.84,0.14) -- (0.84,0.86)
    (0.16,0.14) -- (0.84,0.14) (0.16,0.86) -- (0.84,0.86)
    (0.16,0.60) arc[start angle=180,end angle=0,radius=0.11]
    (0.38,0.60) arc[start angle=180,end angle=0,radius=0.12]
    (0.62,0.60) arc[start angle=180,end angle=0,radius=0.11]
    (0.16,0.74) -- (0.84,0.74);}

% Lupa (investigação)
\newcommand{\icoLupa}[1]{%
  \draw[color=#1,line width=1.1pt] (0.28,0.60) circle[radius=0.20]
    (0.42,0.46) -- (0.70,0.18);
  \draw[color=#1,line width=1.1pt] (0.21,0.67) arc[start angle=180,end angle=270,radius=0.07];}

% Cana-de-açúcar
\newcommand{\icoCana}[1]{%
  \draw[color=#1,line width=1.1pt,line join=round]
    (0.50,0.10) -- (0.50,0.88)
    (0.50,0.88) -- (0.80,0.68) -- (0.63,0.58)
    (0.50,0.88) -- (0.20,0.68) -- (0.37,0.58)
    (0.42,0.40) -- (0.60,0.52) (0.60,0.28) -- (0.42,0.42);}

% Lápide com cruz
\newcommand{\icoLapide}[1]{%
  \draw[color=#1,line width=1.1pt,line join=round]
    (0.28,0.14) -- (0.28,0.62) arc[start angle=180,end angle=0,radius=0.22] -- (0.72,0.14) -- cycle
    (0.50,0.50) -- (0.50,0.72) (0.40,0.62) -- (0.60,0.62);}

% Cruzeiro de pedra
\newcommand{\icoCruz}[1]{%
  \draw[color=#1,line width=1.1pt,line join=round]
    (0.50,0.90) -- (0.50,0.24) (0.34,0.74) -- (0.66,0.74)
    (0.30,0.10) -- (0.70,0.10) (0.36,0.17) -- (0.64,0.17);}

% Azulejo tipo tapete
\newcommand{\icoAzulejo}[1]{%
  \draw[color=#1,line width=1.1pt]
    (0.16,0.16) rectangle (0.84,0.84)
    (0.16,0.16) -- (0.50,0.50) (0.50,0.84) -- (0.84,0.50)
    (0.50,0.16) -- (0.16,0.50) (0.84,0.84) -- (0.50,0.50)
    (0.16,0.50) -- (0.50,0.84) (0.50,0.50) -- (0.84,0.16);}

% Lajotas de barro (piso)
\newcommand{\icoLajota}[1]{%
  \draw[color=#1,line width=1.1pt]
    (0.16,0.22) rectangle (0.84,0.78)
    (0.16,0.41) -- (0.84,0.41) (0.16,0.60) -- (0.84,0.60)
    (0.44,0.22) -- (0.44,0.41) (0.64,0.41) -- (0.64,0.60) (0.44,0.60) -- (0.44,0.78);}

% Pia de cantaria (taça em pedestal)
\newcommand{\icoPia}[1]{%
  \draw[color=#1,line width=1.1pt,line join=round]
    (0.34,0.74) -- (0.66,0.74)
    (0.66,0.74) arc[start angle=0,end angle=-180,radius=0.16]
    (0.50,0.58) -- (0.50,0.34)
    (0.38,0.34) -- (0.62,0.34) -- (0.62,0.24) -- (0.38,0.24) -- cycle
    (0.40,0.79) -- (0.60,0.79);}

% Porta com cercadura de pedra
\newcommand{\icoPorta}[1]{%
  \draw[color=#1,line width=1.1pt,line join=round]
    (0.24,0.14) rectangle (0.76,0.86)
    (0.32,0.14) -- (0.32,0.56) arc[start angle=180,end angle=0,radius=0.18] -- (0.68,0.14) -- cycle;}

% Tijolos de adobe secando ao sol
\newcommand{\icoTijolo}[1]{%
  \draw[color=#1,line width=1.1pt,line join=round]
    (0.18,0.20) rectangle (0.48,0.42)
    (0.52,0.20) rectangle (0.82,0.42)
    (0.35,0.46) rectangle (0.65,0.68)
    (0.14,0.78) -- (0.24,0.78) (0.19,0.72) -- (0.19,0.84)
    (0.76,0.78) -- (0.86,0.78) (0.81,0.72) -- (0.81,0.84);}

% Reboco (parede recebendo espátula)
\newcommand{\icoReboco}[1]{%
  \draw[color=#1,line width=1.1pt,line join=round]
    (0.24,0.16) rectangle (0.76,0.64)
    (0.24,0.30) -- (0.76,0.30) (0.24,0.46) -- (0.76,0.46)
    (0.45,0.64) -- (0.63,0.82)
    (0.63,0.82) -- (0.79,0.66) -- (0.69,0.56) -- (0.53,0.72) -- cycle;}

% Forno de cal com fumaça
\newcommand{\icoCal}[1]{%
  \draw[color=#1,line width=1.1pt,line join=round]
    (0.20,0.16) -- (0.50,0.70) -- (0.80,0.16) -- cycle
    (0.50,0.70) .. controls (0.42,0.80) and (0.58,0.84) .. (0.50,0.94);}

% Gota d'água
\newcommand{\icoGota}[1]{%
  \draw[color=#1,line width=1.1pt,line join=round]
    (0.50,0.88) .. controls (0.36,0.64) and (0.28,0.54) .. (0.28,0.42)
    arc[start angle=180,end angle=360,radius=0.22]
    .. controls (0.72,0.54) and (0.64,0.64) .. (0.50,0.88) -- cycle
    (0.18,0.12) -- (0.82,0.12);}

% Setas de troca
\newcommand{\icoTroca}[1]{%
  \draw[color=#1,line width=1.1pt,line cap=round]
    (0.80,0.50) arc[start angle=0,end angle=150,radius=0.28]
    (0.20,0.50) arc[start angle=180,end angle=330,radius=0.28];
  \draw[color=#1,line width=1.1pt,line join=round]
    (0.66,0.76) -- (0.80,0.50) -- (0.87,0.79)
    (0.34,0.24) -- (0.20,0.50) -- (0.13,0.21);}

% Vento sul (nuvem e rajadas)
\newcommand{\icoVento}[1]{%
  \draw[color=#1,line width=1.1pt,line join=round]
    (0.26,0.64) .. controls (0.26,0.76) and (0.40,0.78) .. (0.46,0.70)
    .. controls (0.50,0.80) and (0.68,0.80) .. (0.70,0.68)
    .. controls (0.82,0.68) and (0.84,0.56) .. (0.74,0.52)
    -- (0.32,0.52) .. controls (0.24,0.52) and (0.22,0.58) .. (0.26,0.64) -- cycle;
  \draw[color=#1,line width=1.1pt,line cap=round]
    (0.20,0.36) -- (0.68,0.36) (0.30,0.24) -- (0.80,0.24);}

% Sino
\newcommand{\icoSino}[1]{%
  \draw[color=#1,line width=1.1pt,line join=round]
    (0.50,0.88) -- (0.50,0.80)
    (0.36,0.76) -- (0.64,0.76)
    (0.36,0.76) .. controls (0.36,0.56) and (0.40,0.44) .. (0.30,0.30)
    (0.64,0.76) .. controls (0.64,0.56) and (0.60,0.44) .. (0.70,0.30)
    (0.30,0.30) -- (0.70,0.30);
  \draw[color=#1,line width=1.1pt,fill=#1] (0.50,0.24) circle[radius=0.04];}

% Tábuas de forro
\newcommand{\icoForro}[1]{%
  \draw[color=#1,line width=1.1pt,line join=round]
    (0.18,0.20) -- (0.82,0.20) (0.18,0.40) -- (0.82,0.40)
    (0.18,0.60) -- (0.82,0.60) (0.18,0.80) -- (0.82,0.80)
    (0.34,0.20) -- (0.34,0.40) (0.60,0.40) -- (0.60,0.60) (0.42,0.60) -- (0.42,0.80);}

% ============================================================
%  \carta — desenha UMA carta completa
%  #1 cor do naipe      #2 nome do naipe (minúsculo, p/ a banda)
%  #3 pictograma naipe  #4 título
%  #5 valor (nº)        #6 ilustração (código TikZ na caixa 0..1)
%  #7 citação           #8 saber (fato)
%  \cartaB: #1 efeito  #2 código  #3 cor da peça  #4 inicial da peça (vazio = sem selo)
% ============================================================
\newcommand{\carta}[8]{%
  \begingroup
  \def\naipecor{#1}\def\naipenome{#2}\def\naipeicon{#3}%
  \def\cartatitulo{#4}\def\cartavalor{#5}%
  \def\cartaarte{#6}\def\cartacitacao{#7}%
  \def\cartasaber{#8}%
  \cartaB
}
\def\vazio{}
\newcommand{\cartaB}[4]{%
  \def\cartaeff{#1}\def\cartacodigo{#2}\def\partecor{#3}\def\parteinicial{#4}%
  \begin{tikzpicture}[x=1mm,y=1mm]
    % --- corpo ---
    \draw[rounded corners=2.5mm,fill=white,draw=black!45,line width=0.4pt]
      (0,0) rectangle (\larguracarta,\alturacarta);
    % --- banda lateral esquerda ---
    \fill[\naipecor,rounded corners=2.5mm]
      (0,0) rectangle (7.5mm,\alturacarta);
    \fill[\naipecor]
      (3.75mm,0) rectangle (7.5mm,\alturacarta);
    \node[rotate=90,font=\sffamily\bfseries\footnotesize,text=white]
      at (3.75mm,{0.5\alturacarta}) {\MakeUppercase{\naipenome}};
    % --- faixa superior ---
    \fill[\naipecor,rounded corners=1.2mm] (9.5mm,79.5mm) rectangle (17.5mm,87.5mm);
    \begin{scope}[shift={(10.1mm,80.1mm)},x=7.4mm,y=7.4mm]
      \naipeicon
    \end{scope}
    \fill[tinta,rounded corners=1.2mm] (19mm,79.5mm) rectangle (52.5mm,87.5mm);
    \node[font=\sffamily\bfseries\footnotesize,text=white,text centered,text width=32mm]
      at (35.75mm,83.5mm) {\cartatitulo};
    \fill[ouro,rounded corners=1.2mm] (54mm,79.5mm) rectangle (62mm,87.5mm);
    \node[font=\sffamily\bfseries\large,text=tinta] at (58mm,83.5mm) {\cartavalor};
    % --- ilustração ---
    \fill[areia,rounded corners=1.5mm] (9.5mm,48.5mm) rectangle (62mm,78mm);
    \begin{scope}[shift={(11.5mm,49.7mm)},x=48.5mm,y=26.6mm]
      \cartaarte
    \end{scope}
    % --- citação ---
    \node[anchor=north west,font=\itshape\scriptsize,text=tinta,text width=50mm,inner sep=0pt]
      at (9.5mm,46.8mm) {``\cartacitacao''};
    \fill[tinta] (9.5mm,39.8mm) rectangle (62mm,40.6mm);
    % --- saber (label em linha com o texto) ---
    \node[anchor=north west,font=\fontsize{6.5}{7.8}\selectfont,text=tinta,text width=51mm,inner sep=0pt]
      at (9.5mm,38.2mm) {{\color{\naipecor}\sffamily\bfseries SABER\ \ }\cartasaber};
    % --- efeito (label em linha com o texto) ---
    \node[anchor=north west,font=\fontsize{6.5}{7.8}\selectfont,text=tinta,text width=51mm,inner sep=0pt]
      at (9.5mm,22.6mm) {{\color{black!60}\sffamily\bfseries EFEITO\ \ }\cartaeff};
    % --- rodapé ---
    \node[anchor=south west,font=\tiny\sffamily,text=black!55,inner sep=0pt]
      at (9.5mm,2mm) {Restauro de Passé};
    \ifx\parteinicial\vazio
      % cartas de ação: sem selo de peça
    \else
      \node[anchor=south east,circle,fill=\partecor,minimum size=4.2mm,inner sep=0pt,
            font=\sffamily\bfseries\scriptsize,text=white]
        at (62mm,4.2mm) {\parteinicial};
    \fi
    \node[anchor=north east,font=\tiny\sffamily\bfseries,text=black!55,inner sep=0pt]
      at (62mm,11mm) {\cartacodigo};
  \end{tikzpicture}%
  \endgroup
}

% ---------- Verso da carta ----------
\newcommand{\verso}{%
  \begin{tikzpicture}[x=1mm,y=1mm]
    \draw[rounded corners=2.5mm,fill=historia,draw=black!45,line width=0.4pt]
      (0,0) rectangle (\larguracarta,\alturacarta);
    \draw[rounded corners=2mm,fill=white,draw=white] (4mm,4mm) rectangle (59mm,84mm);
    \begin{scope}[shift={(6mm,14mm)},x=51mm,y=60mm]
      \icoIgreja{historia}
    \end{scope}
    \node[font=\sffamily\bfseries\small,text=historia,text centered,text width=50mm]
      at (31.5mm,8mm) {RESTAURO DE PASSÉ};
  \end{tikzpicture}%
}

% ---------- Helpers de composição ----------
% Uma linha de 3 cartas (use \par ou linha em branco entre linhas de cartas)
\newcommand{\linhacartas}[3]{#1\hspace{2.2mm}#2\hspace{2.2mm}#3\par}
% Uma linha de 3 versos
\newcommand{\linhaversos}{\linhacartas{\verso}{\verso}{\verso}}

\setlength{\parindent}{0pt}
\pagestyle{empty}


\begin{document}

% ==================== PÁGINA 1 — HISTÓRIA (1/2) ====================
\linhacartas
{\carta{historia}{história}{\icoPergaminho{white}}{A PORTA DE PEDRA}{1}{\icoPorta{tinta}}
{Por aqui passaram gerações de passenses.}
{A entrada principal é em pedra de cantaria: blocos cortados à mão por mestres pedreiros, com cercadura em volta. Pena que o tempo levou a madeira da porta.}
{Guarde: peça Fachada e torre.}{H01}{fachada}{F}}
{\carta{historia}{história}{\icoPergaminho{white}}{TORRES SEM FIM}{2}{\icoTorre{tinta}}
{Um dia as torres vão ganhar sinos\ldots}
{As duas torres da fachada ficaram inacabadas --- e assim seguem até hoje, guardando o céu de Passé como duas sentinelas de pedra.}
{Guarde: peça Fachada e torre.}{H02}{fachada}{F}}
{\carta{historia}{história}{\icoPergaminho{white}}{O FRONTÃO}{2}{\icoIgreja{tinta}}
{O rosto da igreja mudou no século XVIII.}
{Na reforma da segunda metade do séc. XVIII, a fachada ganhou cornija curva e frontão recortado, detalhes que ainda enfeitam as ruínas.}
{Guarde: peça Fachada e torre.}{H03}{fachada}{F}}
\linhacartas
{\carta{historia}{história}{\icoPergaminho{white}}{O ADRO MAIS BONITO}{1}{\icoBarco{tinta}}
{Lá embaixo, a Baía de Todos-os-Santos.}
{A igreja foi erguida no alto de uma elevação que cai sobre a baía. Do adro, debruçado no muro, avista-se a Ilha de Maré em primeiro plano.}
{Guarde: peça Fachada e torre.}{H04}{fachada}{F}}
{\carta{historia}{história}{\icoPergaminho{white}}{A GALERIA DUPLA}{3}{\icoGaleria{tinta}}
{Corredores de sombra e vento.}
{A nave tinha galeria dupla aberta para o exterior --- e especialistas dizem que talvez seja a mais antiga igreja baiana com essa raridade.}
{Guarde: peça Nave.}{H05}{nave}{N}}
{\carta{historia}{história}{\icoPergaminho{white}}{A NAVE DO POVO}{1}{\icoIgreja{tinta}}
{Aqui a freguesia se encontrava.}
{A nave é o salão principal da igreja, onde os moradores rezavam, se encontravam e celebravam a vida do povoado.}
{Guarde: peça Nave.}{H06}{nave}{N}}
\linhacartas
{\carta{historia}{história}{\icoPergaminho{white}}{298 CASAS}{2}{\icoCasas{tinta}}
{Uma freguesia quase uma cidade!}
{Em 1759, um registro contou 298 casas e 2.497 habitantes na freguesia de Nossa Senhora do Passé --- povoação de gente trabalhadora.}
{Guarde: peça Nave.}{H07}{nave}{N}}
{\carta{historia}{história}{\icoPergaminho{white}}{A SESMARIA}{2}{\icoPergaminho{tinta}}
{Tudo começou com uma carta de terra.}
{Entre 1563 e 1566, jesuítas receberam a sesmaria de Passé, com grandes engenhos de cana como Pitanga, Jacarenga e Restinga.}
{Guarde: peça Nave.}{H08}{nave}{N}}
{\carta{historia}{história}{\icoPergaminho{white}}{O MISTÉRIO DAS DATAS}{3}{\icoLupa{tinta}}
{1620? Fim do séc. XVII? Ainda se estuda!}
{A data da construção divide os estudiosos: a prefeitura já divulgou 1620, mas o inventário do IPAC aponta o fim do séc. XVII. Talvez você ajude a resolver!}
{Guarde: peça Nave.}{H09}{nave}{N}}
\newpage

% ==================== PÁGINA 2 — HISTÓRIA (2/2) ====================
\linhacartas
{\carta{historia}{história}{\icoPergaminho{white}}{CANA E CORAGEM}{2}{\icoCana{tinta}}
{A doçura do açúcar tinha um preço triste.}
{A riqueza da freguesia veio da cana-de-açúcar, produzida por africanos e indígenas escravizados. Essa memória também merece ser contada.}
{Guarde: peça Nave.}{H10}{nave}{N}}
{\carta{historia}{história}{\icoPergaminho{white}}{O CORAÇÃO DA IGREJA}{2}{\icoIgreja{tinta}}
{Perto do altar, tudo era mais sagrado.}
{A capela-mor é a parte mais nobre da igreja, onde ficavam o altar e as imagens, entre elas a de Nossa Senhora da Encarnação.}
{Guarde: peça Capela-mor.}{H11}{capela}{C}}
{\carta{historia}{história}{\icoPergaminho{white}}{SEPULTADOS NA IGREJA}{1}{\icoLapide{tinta}}
{Quanto mais perto do altar, mais perto do céu.}
{Antigamente enterrava-se dentro das igrejas, e o lugar junto ao altar era o mais caro --- só para quem podia pagar.}
{Guarde: peça Capela-mor.}{H12}{capela}{C}}
\linhacartas
{\carta{historia}{história}{\icoPergaminho{white}}{LÁPIDES QUE FALAM}{3}{\icoLapide{tinta}}
{Sempre chorada, nunca esquecida.}
{No chão da igreja há lápides de 1883 e de 1932, como a de Maria Bertholina Gomes Betuca. Pedras que guardam nomes e saudades.}
{Guarde: peça Capela-mor.}{H13}{capela}{C}}
{\carta{historia}{história}{\icoPergaminho{white}}{FESTA DE 25 DE MARÇO}{1}{\icoSino{tinta}}
{A padroeira reúne o povoado até hoje.}
{Devotos celebram Nossa Senhora da Encarnação em 25 de março, na localidade de Mucunga, distrito de Passé --- a fé continua viva.}
{Guarde: peça Capela-mor.}{H14}{capela}{C}}
{\carta{historia}{história}{\icoPergaminho{white}}{PAREDES MISTERIOSAS}{1}{\icoPilao{tinta}}
{Será taipa? Será adobe? Ainda se investiga!}
{Os registros históricos não contam como as paredes foram erguidas. Arquitetos e historiadores seguem estudando as ruínas para descobrir.}
{Guarde: peça Paredes.}{H15}{taipa}{T}}
\linhacartas
{\carta{historia}{história}{\icoPergaminho{white}}{UM TELHADO SÓ}{1}{\icoTelha{tinta}}
{Duas águas para uma igreja inteira.}
{A igreja tinha telhado único de duas águas cobrindo nave e capela-mor: a chuva escorria para os lados, longe das paredes.}
{Guarde: peça Telhado.}{H16}{telhado}{S}}
{\carta{construcao}{construção}{\icoTaipa{white}}{CERCADURAS E CUNHAIS}{2}{\icoArco{tinta}}
{A pedra firme nas quinas e nas janelas.}
{A cantaria de pedra reforçava portada, janelas, cantos (cunhais) e o arco-cruzeiro --- justamente as partes que mais peso carregam.}
{Guarde: peça Fachada e torre.}{C01}{fachada}{F}}
{\carta{construcao}{construção}{\icoTaipa{white}}{AZULEJOS TIPO TAPETE}{3}{\icoAzulejo{tinta}}
{Azul e amarelo como o céu e o sol.}
{Sob o coro havia azulejos portugueses seiscentistas, ``tipo tapete'': desenhos azuis e amarelos sobre o branco, colados um a um.}
{Guarde: peça Nave.}{C02}{nave}{N}}
\newpage

% ==================== PÁGINA 3 — CONSTRUÇÃO (1/2) ====================
\linhacartas
{\carta{construcao}{construção}{\icoTaipa{white}}{LAJOTAS NO CHÃO}{1}{\icoLajota{tinta}}
{Um chão fresco feito de barro.}
{Nave e sacristia tinham piso de lajotas de barro --- peças grandes assentadas uma a uma, frescas no calor do Recôncavo.}
{Guarde: peça Nave.}{C03}{nave}{N}}
{\carta{construcao}{construção}{\icoTaipa{white}}{A PIA BATISMAL}{1}{\icoPia{tinta}}
{Nela começou a história de muitos passenses.}
{O inventário registrou pias de cantaria na igreja: a de água benta, na entrada, e a batismal, onde nasciam de novo os fiéis.}
{Guarde: peça Nave.}{C04}{nave}{N}}
{\carta{construcao}{construção}{\icoTaipa{white}}{ALPENDRE DE SOMBRA}{2}{\icoGaleria{tinta}}
{Chuva de fora, frescor de dentro.}
{O alpendre lateral protegia as entradas da chuva e do sol forte --- uma técnica simples que ainda vale para as casas de hoje.}
{Guarde: peça Nave.}{C05}{nave}{N}}
\linhacartas
{\carta{construcao}{construção}{\icoTaipa{white}}{O CORO DE MADEIRA}{2}{\icoForro{tinta}}
{De cima, o canto guiava a rezada.}
{O coro, sobre a entrada, era de madeira; embaixo dele ficavam os azulejos que enfeitavam a nave --- som e beleza no mesmo lugar.}
{Guarde: peça Nave.}{C06}{nave}{N}}
{\carta{construcao}{construção}{\icoTaipa{white}}{ABÓBADA DE MADEIRA}{2}{\icoForro{tinta}}
{Madeira fingindo pedra.}
{O forro da capela-mor imitava uma abóbada abatida, toda em madeira moldada --- carpintaria fina sem usar uma pedra sequer.}
{Guarde: peça Capela-mor.}{C07}{capela}{C}}
{\carta{construcao}{construção}{\icoTaipa{white}}{PÚLPITO E LAVABO}{1}{\icoPia{tinta}}
{Cada pedra tinha um ofício.}
{A cantaria servia a tudo: bacia do púlpito para pregar, lavabo para as cerimônias e embasamento para segurar a capela-mor.}
{Guarde: peça Capela-mor.}{C08}{capela}{C}}
\linhacartas
{\carta{construcao}{construção}{\icoTaipa{white}}{TAIPA DE PILÃO}{2}{\icoTaipa{tinta}}
{Terra socada que vira parede.}
{Na taipa de pilão, a terra úmida é socada dentro de um molde de madeira, em camadas de uns 40 cm. Quando seca, o molde sobe --- parede sem tijolo nem cimento!}
{Guarde: peça Paredes.}{C09}{taipa}{T}}
{\carta{construcao}{construção}{\icoTaipa{white}}{ADOBE AO SOL}{1}{\icoTijolo{tinta}}
{Tijolo de barro que o sol cozinha.}
{O adobe é feito de barro moldado e secado ao sol, sem queima. É leve, fresco, barato e fácil de consertar --- perfeito para mutirão.}
{Guarde: peça Paredes.}{C10}{taipa}{T}}
{\carta{construcao}{construção}{\icoTaipa{white}}{O REBOCO PROTETOR}{1}{\icoReboco{tinta}}
{A pele da parede contra a chuva.}
{O reboco de cal e areia cobre e protege a parede de terra, levando a água embora. Sem reboco, a taipa se dissolve na primeira chuva.}
{Guarde: peça Paredes.}{C11}{taipa}{T}}
\newpage

% ==================== PÁGINA 4 — CONSTRUÇÃO (2/2) + DESAFIO ====================
\linhacartas
{\carta{construcao}{construção}{\icoTaipa{white}}{PEDRA QUE VIRA COLA}{1}{\icoCal{tinta}}
{A cal nasce da pedra queimada.}
{A cal é feita queimando pedra calcária; misturada com areia e água vira argamassa, que cola as pedras e recebe os rebocos.}
{Guarde: peça Paredes.}{C12}{taipa}{T}}
{\carta{construcao}{construção}{\icoTaipa{white}}{TELHA VÃ}{2}{\icoTelha{tinta}}
{Na nave, o telhado aparecia sem vergonha.}
{A nave usava ``telha vã'': o telhado ficava à mostra, sem forro escondido --- bonito, arejado e fácil de vistoriar.}
{Guarde: peça Telhado.}{C13}{telhado}{S}}
{\carta{construcao}{construção}{\icoTaipa{white}}{CANAL E CAPA}{1}{\icoGota{tinta}}
{Um telhado de barro é um telhado de casal.}
{A telha colonial é dupla: a canal escorre a água e a capa a protege. São peças de barro moldadas uma a uma, na coxa ou em fôrma.}
{Guarde: peça Telhado.}{C14}{telhado}{S}}
\linhacartas
{\carta{construcao}{construção}{\icoTaipa{white}}{MADEIRAMENTO}{2}{\icoForro{tinta}}
{Tesoura, caibro e ripa segurando o céu.}
{O telhado pousa num esqueleto de madeira: tesouras grandes, caibros e ripas --- quase sempre de árvores da própria região.}
{Guarde: peça Telhado.}{C15}{telhado}{S}}
{\carta{construcao}{construção}{\icoTaipa{white}}{VERGA E CONTRA-VERGA}{1}{\icoArco{tinta}}
{Sobre a porta, a madeira segura o mundo.}
{Portas e janelas recebem verga e contra-verga para carregar o peso da parede por cima sem rachar --- detalhe que toda casa precisa.}
{Guarde: peça Telhado.}{C16}{telhado}{S}}
{\carta{desafio}{desafio}{\icoDesafio{white}}{QUE SÉCULO FOI?}{2}{\icoLupa{tinta}}
{A data é o maior mistério do monumento.}
{Pergunta: segundo o inventário do IPAC, de quando é a igreja --- começo do séc. XVII, fim do séc. XVII ou séc. XIX?\\ \textbf{Resposta:} fim do séc. XVII. (A prefeitura já divulgou 1620.)}
{Acertou? Fique com a carta (curinga). Errou? Fundo do monte.}{D01}{ouro}{$\star$}}
\linhacartas
{\carta{desafio}{desafio}{\icoDesafio{white}}{QUEM COMEÇOU TUDO?}{2}{\icoPergaminho{tinta}}
{Toda história tem um primeiro capítulo.}
{Pergunta: quem recebeu a sesmaria de Passé entre 1563 e 1566?\\ \textbf{Resposta:} os jesuítas, donos de grandes engenhos de cana.}
{Acertou? Fique com a carta (curinga). Errou? Fundo do monte.}{D02}{ouro}{$\star$}}
{\carta{desafio}{desafio}{\icoDesafio{white}}{A FREGUESIA CRESCIA}{2}{\icoCasas{tinta}}
{Passé já foi quase uma cidade.}
{Pergunta: quantas casas a freguesia tinha em 1759?\\ \textbf{Resposta:} 298 casas, com 2.497 habitantes.}
{Acertou? Fique com a carta (curinga). Errou? Fundo do monte.}{D03}{ouro}{$\star$}}
{\carta{desafio}{desafio}{\icoDesafio{white}}{TORRES OU MIRANTES?}{2}{\icoTorre{tinta}}
{Olhe a fachada: falta alguma coisa.}
{Pergunta: o que dizem as duas torres da igreja?\\ \textbf{Resposta:} que elas nunca foram terminadas --- estão inacabadas até hoje.}
{Acertou? Fique com a carta (curinga). Errou? Fundo do monte.}{D04}{ouro}{$\star$}}
\newpage

% ==================== PÁGINA 5 — DESAFIO (2/2) ====================
\linhacartas
{\carta{desafio}{desafio}{\icoDesafio{white}}{O QUE É TAIPA?}{2}{\icoTaipa{tinta}}
{Técnica de parede de terra socada.}
{Pergunta: como se levanta uma parede de taipa de pilão?\\ \textbf{Resposta:} socando terra úmida dentro de um molde de madeira, em camadas.}
{Acertou? Fique com a carta (curinga). Errou? Fundo do monte.}{D05}{ouro}{$\star$}}
{\carta{desafio}{desafio}{\icoDesafio{white}}{TELHA SEM FORRO}{2}{\icoTelha{tinta}}
{A nave tinha o teto à mostra.}
{Pergunta: o que os antigos chamavam de ``telha vã''?\\ \textbf{Resposta:} telhado aparente, sem forro escondido.}
{Acertou? Fique com a carta (curinga). Errou? Fundo do monte.}{D06}{ouro}{$\star$}}
{\carta{desafio}{desafio}{\icoDesafio{white}}{ILHA NO HORIZONTE}{2}{\icoBarco{tinta}}
{Do adro, o olhar encontra a baía.}
{Pergunta: que ilha se vê da frente da igreja?\\ \textbf{Resposta:} a Ilha de Maré, na Baía de Todos-os-Santos.}
{Acertou? Fique com a carta (curinga). Errou? Fundo do monte.}{D07}{ouro}{$\star$}}
\linhacartas
{\carta{desafio}{desafio}{\icoDesafio{white}}{DIA DE FESTA}{2}{\icoSino{tinta}}
{A devoção atravessa os séculos.}
{Pergunta: quando é a festa de Nossa Senhora da Encarnação em Mucunga?\\ \textbf{Resposta:} em 25 de março.}
{Acertou? Fique com a carta (curinga). Errou? Fundo do monte.}{D08}{ouro}{$\star$}}
{\carta{desafio}{desafio}{\icoDesafio{white}}{TIJOLO SEM FOGO}{2}{\icoTijolo{tinta}}
{Um tijolo mais fresco que o comum.}
{Pergunta: o que é o adobe?\\ \textbf{Resposta:} tijolo de barro moldado e secado ao sol, sem queima.}
{Acertou? Fique com a carta (curinga). Errou? Fundo do monte.}{D09}{ouro}{$\star$}}
{\carta{desafio}{desafio}{\icoDesafio{white}}{PÉS NO BARRO}{2}{\icoLajota{tinta}}
{O chão da nave ainda conta sua história.}
{Pergunta: com que material era o piso da nave e da sacristia?\\ \textbf{Resposta:} lajotas de barro.}
{Acertou? Fique com a carta (curinga). Errou? Fundo do monte.}{D10}{ouro}{$\star$}}
\linhacartas
{\carta{desafio}{desafio}{\icoDesafio{white}}{ESCUDO CONTRA A CHUVA}{2}{\icoReboco{tinta}}
{Toda parede de terra precisa de proteção.}
{Pergunta: o que protege a taipa e a pedra da chuva?\\ \textbf{Resposta:} o reboco de cal e areia.}
{Acertou? Fique com a carta (curinga). Errou? Fundo do monte.}{D11}{ouro}{$\star$}}
{\carta{desafio}{desafio}{\icoDesafio{white}}{INVENTÁRIO NÃO É TOMBAMENTO}{2}{\icoLupa{tinta}}
{Em 1977 o IPAC fichou a igreja --- mas falta um passo.}
{Pergunta: ser inventariada significa ser tombada?\\ \textbf{Resposta:} não! O tombamento, que protege de verdade, ainda é conquista da comunidade.}
{Acertou? Fique com a carta (curinga). Errou? Fundo do monte.}{D12}{ouro}{$\star$}}
{\carta{acao}{ação}{\icoAcao{white}}{PILÃO! PILÃO!}{1}{\icoPilao{tinta}}
{Um enche, outro soca, todos cantam.}
{Erguer taipa é trabalho de mutirão: a obra rende porque a comunidade trabalha junta, na cadência do pilão.}
{Compre 2 cartas. Leia os dois SABERes em voz alta e fique com 1; a outra volta para o fundo do monte.}{A01}{}{}}
\newpage

% ==================== PÁGINA 6 — AÇÃO ====================
\linhacartas
{\carta{acao}{ação}{\icoAcao{white}}{TROCA-TROCA}{1}{\icoTroca{tinta}}
{Meu tijolo pelo seu favor.}
{Nas obras antigas, trocar materiais e mãos amigas com os vizinhos era o jeito de construir --- e de fazer comunidade.}
{Troque 1 carta de peça com outro jogador. Cada um escolhe a sua.}{A02}{}{}}
{\carta{acao}{ação}{\icoAcao{white}}{MESTRE DE OBRAS}{2}{\icoAcao{tinta}}
{A obra inteira de cor e salteado.}
{O mestre de obras comandava tudo sem desenho na mão: sabia de olho onde cada pedra e cada viga iam parar.}
{Diga uma peça que falta para você. Cada jogador que tiver carta dessa peça deve entregar 1 a você.}{A03}{}{}}
{\carta{acao}{ação}{\icoAcao{white}}{CHUVA FORTE}{1}{\icoGota{tinta}}
{Trancá tudo e reza, que lá vem água!}
{A chuva é a maior inimiga das paredes de terra e do reboco de cal. É por isso que o telhado e o alpendre existem.}
{Perdeu a vez! Para a chuva passar mais rápido, leia o SABER em voz alta.}{A04}{}{}}
\linhacartas
{\carta{acao}{ação}{\icoAcao{white}}{ROMARIA}{1}{\icoSino{tinta}}
{Fé, festa e quitutes.}
{As romarias reuniam povos de várias freguesias na igrejinha: muita festa, muita comida e canto até cansar.}
{Todos gritam: ``Viva Nossa Senhora da Encarnação!'' Depois você compra 3 cartas e fica com 1.}{A05}{}{}}
{\carta{acao}{ação}{\icoAcao{white}}{TESOURO DO CORO}{2}{\icoAzulejo{tinta}}
{Que brilho escondido embaixo do coro!}
{Sob o coro escondiam-se azulejos seiscentistas ``tipo tapete'', azuis e amarelos sobre o branco.}
{Vire cartas do monte até achar uma DESAFIO. Leia o SABER, fique com ela e mande as demais para o fundo do monte.}{A06}{}{}}
{\carta{acao}{ação}{\icoAcao{white}}{VENTO SUL}{1}{\icoVento{tinta}}
{Segura a telha que o vento quer levar!}
{Vento forte levanta telha solta --- por isso a capa pesa sobre a canal e a ripa prende tudo no lugar.}
{Todos os jogadores passam 1 carta Telhado para a pessoa à sua esquerda. Quem não tiver, passa a vez de passar.}{A07}{}{}}
\linhacartas
{\carta{acao}{ação}{\icoAcao{white}}{RESTAURO DO IPAC}{2}{\icoIgreja{tinta}}
{Consertar com carinho, sem apagar o tempo.}
{Restaurar um bem histórico é cuidar sem mentir: as marcas do tempo ficam à mostra, e a comunidade participa de tudo.}
{Troque 1 carta da sua frente por 1 carta do topo do monte. Leia o SABER da nova carta em voz alta.}{A08}{}{}}
{\carta{acao}{ação}{\icoAcao{white}}{FESTA DA PADROEIRA}{1}{\icoSol{tinta}}
{25 de março é dia de alegria!}
{A festa de Nossa Senhora da Encarnação mantém viva a devoção em Mucunga, no distrito de Passé --- a tradição atravessa gerações.}
{Cada jogador lê 1 SABER das suas cartas em voz alta. Depois você compra 2 e fica com 1.}{A09}{}{}}
{\carta{acao}{ação}{\icoAcao{white}}{RISADA NO TERREIRO}{1}{\icoAcao{tinta}}
{O pilão soa e a obra vira festa.}
{A obra era também alegria: canto, conversa e comida dividida no terreiro, ao som ritmado do pilão.}
{Imite o barulho do pilão de taipa! Se alguém rir, ganhe 1 carta de peça dessa pessoa --- ela escolhe qual.}{A10}{}{}}
\newpage

% ==================== VERSOS (imprimir no verso das folhas) ====================
\linhaversos
\linhaversos
\linhaversos
\newpage
\linhaversos
\linhaversos
\linhaversos
\newpage
\linhaversos
\linhaversos
\linhaversos
\newpage
\linhaversos
\linhaversos
\linhaversos
\newpage
\linhaversos
\linhaversos
\linhaversos
\newpage
\linhaversos
\linhaversos
\linhaversos

\end{document}
